\subsection{A separate unused Gaia astrometric test}
\label{sec:he0435-gaia}
The LCO common-reference limitation motivates a second type of observation:
the four image positions in Gaia DR3~\cite{GaiaDR3}. We declare this
additional test after the LCO result is known, while retaining exactly the
already-frozen conventional and R1 physical predictions. A count-only
catalogue query establishes four entries in the lens field. Numerical
positions, errors and quality rows remain unopened until a separate
prediction/protocol commit. No mass, screen, source, angle, plate scale or
rotation is fitted against Gaia.

In a common small-angle ICRS tangent frame the three relative observables
are $z_i=((\alpha_i-\alpha_A)\cos\delta_A,\delta_i-\delta_A)$,
$i=B,C,D$, with degree-to-arcsecond conversion. All 24 source assignments
are compared with the fixed conventional relative pattern; the same
assignment is used for R1 and its runner-up gap is retained. For each Gaia
source, the published RA tangent-plane error
$\sigma_{\alpha*}=\sigma_\alpha\cos\delta$, declination error and their
correlation define $C_i$. The RA error receives no second cosine.
The full six-coordinate formal covariance is
\begin{equation}
 [C_{\rm rel}]_{ij}=\delta_{ij}C_i+C_A,
 \qquad C_s=C_{\rm rel}+s^2I_6,
 \quad s=0,2,5\ {\rm mas}.
 \label{eq:gaia-covariance}
\end{equation}
The added floors apply per relative coordinate after differencing; they
are prespecified sensitivity controls, not fitted noise or a derived macro
posterior. The Gaussian density and quadratic residual use $C_s$, including
its determinant and common-A off-diagonal blocks. Between-source catalogue
covariances are unpublished, so source-error independence remains a
conditional approximation. All four sources are retained irrespective of
their quality diagnostics; missing mandatory errors cannot be replaced
silently with zero correlation.

Gaia collected different photons in 2014--2017, whereas the fitted NIC2
point astrometry was published in 2011. However, the reprocessed HST central
mosaics use a Gaia-EDR3 field frame. Their pixel test uses relative
registration, but independence of every field calibration is not certified.
Close-image blending, chromatic photocentres and fixed macro parameters
further limit a catalogue-only Gaussian interpretation. This test assesses
the previously unused numerical quasar geometry; it is not a complete
Einstein-ring reconstruction or a BHSM state measurement.

The Gaia prediction freeze was committed before the numerical query.
The unique closest A/B/C/D assignment has squared residual
$6.3594\times10^{-5}$ arcsec$^2$; the runner-up is $4.6197$ arcsec$^2$.
No continuous frame parameter is fitted. The conventional observed-minus-
predicted east/north residuals are $(0.623,1.701)$ mas for B,
$(1.915,-5.026)$ mas for C and $(-5.572,0.580)$ mas for D.
Their coordinate RMS is $3.25560$ mas, versus $3.25557$ mas for R1.
Table~\ref{tab:gaia-frozen} reports every prespecified covariance control.

\begin{table}[htbp]
 \centering
 \caption{Previously unused Gaia DR3 relative astrometry, with all physical
 predictions fixed before the query. $Q=r^TC_s^{-1}r$ uses six coordinates;
 the log density is normalized in mas$^{-6}$. Added floors are sensitivity
 controls, not fitted errors or a completed macro posterior.}
 \label{tab:gaia-frozen}
 \begin{tabular}{lrrrr}
 \toprule
 Floor (mas) & Conventional $Q$ & R1 $Q$ & Conventional $\log p$ & R1 $\log p$\\
 \midrule
 $0$ & $3465.905$ & $3465.667$ & $-1727.177$ & $-1727.058$\\
 $2$ & $15.7966$ & $15.7963$ & $-17.5910$ & $-17.5908$\\
 $5$ & $2.54112$ & $2.54107$ & $-16.4441$ & $-16.4440$\\
 \bottomrule
 \end{tabular}
\end{table}

Both frozen point predictions fail the catalogue-error-only precision
check. The source RUWE values are $1.00$--$1.11$, none is marked duplicated,
and multi-peak fractions are 0--1\%. B and D nevertheless have odd-window
fractions of 64\% and 69\%, and C has excess-noise significance 2.57;
these diagnostics remain visible without post-target image exclusion.
Adding the declared floors reduces the
formal mismatch, but cannot turn a sensitivity calculation into successful
validation or establish the required physical uncertainty. The R1-minus-
conventional log densities are $0.1191$, $1.53\times10^{-4}$ and
$2.43\times10^{-5}$ for the three controls. A small difference between
inadequate point predictions is not evidence for a BHSM state.
The next physical inference requires justified training/frame/photocentre
and macro-model uncertainty together with richer mass/environment
constraints. These Gaia targets are now consumed; fitting them later
cannot count as another unused-data success.

